\par\noindent\begin{minipage}{\linewidth}
% !TeX root = ../教材题库_全部条目_转写工作稿.tex
% Source: XJ-B1-C02-HW-P044-Y1102
% 第2章 / 2.1 相等关系与不等关系 / 教材 p.44 / PDF 第51页
% 已对照原页转写题干及公式；待教师复核；未制作答案
% 原题及相邻原解截图（核对用）：../题目原图/XJ-B1-C02-HW-P044-Y1102-01.webp
% 本文件只在转写完成后作为题干使用；OCR 条目尚不是可靠题干。
\begin{bankitem}{习题2.1 第11题}{第2章 / 2.1 相等关系与不等关系 / 教材 p.44 / PDF 第51页}{XJ-B1-C02-HW-P044-Y1102}{已对照原页转写题干及公式；待教师复核；未制作答案}
某商品计划提价两次，有甲、乙、丙三种方案，其中 $p>q>0$。经两次提价后，哪种方案提价的幅度大？为什么？
\begin{center}\begin{tabular}{c|c|c}\hline
方案&第一次提价&第二次提价\\\hline
甲&$p\%$&$q\%$\\
乙&$q\%$&$p\%$\\
丙&$\dfrac{p+q}{2}\%$&$\dfrac{p+q}{2}\%$\\\hline
\end{tabular}\end{center}
\end{bankitem}
\end{minipage}\par\medskip
